\section{Related Work}
\label{related_work}

\subsection{Requirement Validation}
\label{rw:validation}

Requirement validation checks whether requirements are correct, complete, consistent, and acceptable to stakeholders.
This paper narrows the validation scope to feasibility and realism: whether a requirement can plausibly be implemented and deployed under known constraints.

\subsection{Requirement Feasibility Analysis}
\label{rw:feasibility}

Feasibility analysis studies whether requirements can be satisfied given technical, organizational, economic, and operational limits.
Our work treats feasibility as an explicit verification target rather than as an informal review activity.

\subsection{Constraint Analysis in Requirements Engineering}
\label{rw:constraint}

Constraint analysis provides a basis for identifying limits related to resources, performance, deployment, and time.
The proposed study uses constraint categories as the organizing structure for feasibility reasoning.

\subsection{LLM-Based Requirements Engineering}
\label{rw:llm_re}

LLMs can help generate, rewrite, and inspect requirements, but generated requirements may be fluent without being realistic.
We position LLM-based agents as feasibility reviewers whose outputs must be grounded in knowledge and validated by humans.

\subsection{Knowledge-Guided Requirement Analysis}
\label{rw:knowledge}

Knowledge-guided analysis incorporates domain knowledge, requirement patterns, constraint rules, and project context into requirement inspection.
This direction is closely related to controlled and structured requirement specification efforts~\cite{darifControlledNaturalLanguage2026}.

\subsection{Multi-Agent Collaboration for Requirement Verification}
\label{rw:multi_agent}

Multi-agent collaboration can separate interpretation, constraint checking, and feasibility judgment into specialized roles.
In this paper, agent collaboration is used as a practical workflow for requirement verification, not as a general-purpose software development framework.
